\section{Representative Material Parameter Sets}

The values below reproduce particular numerical models from papers and
repository implementations authored or coauthored by Zhifeng Zhu.  They are
model-specific parameter sets rather than universal material constants.  In
particular, quantities associated with a heavy-metal layer, spin polarization,
device geometry, or numerical integration are identified separately from
intrinsic magnetic parameters.  Fixed-length unit vectors are used in all three
source-model families summarized here.

\subsection{Property of Heavy Metal}

The spin-Hall angle belongs to the heavy-metal spin-current source and is not
specific to a ferromagnetic, ferrimagnetic, or antiferromagnetic target layer.
For the signed representative values in Table~\ref{tab:heavy-metal-properties},
we use
$\vect{J}_{\mathrm{s}}=\theta_{\mathrm{SH}}\,
\vect{\sigma}\times\vect{J}_{\mathrm{c}}$ and choose Pt as the positive
reference; Ta and W then have negative signs~\cite{sot-stt-compact-model,
liu2012ta-spin-hall,pai2012w-spin-hall}.  Reversing the current, layer-normal,
or spin-polarization convention requires a corresponding sign change in the
torque implementation.

\begin{table}[htbp]
    \centering
    \small
    \renewcommand{\arraystretch}{1.15}
    \begin{tabular}{p{0.16\linewidth}p{0.31\linewidth}p{0.42\linewidth}}
        \hline
        Property & Value & Evidence and qualification \\
        \hline
        Pt spin-Hall angle
            & $+0.07$
            & Representative literature value quoted by Liu et al. and Pai et
              al.; defines the positive reference sign \\
        High-resistivity $\beta$-Ta spin-Hall angle
            & $-0.12$ to $-0.15$
            & Consistent magnitudes from ST-FMR, perpendicular switching, and
              three-terminal switching \\
        High-resistivity $\beta$-W spin-Hall angle
            & $-0.30\pm0.02$ (ST-FMR);
              $-0.33\pm0.06$ (switching)
            & Negative sign verified relative to Pt; the ST-FMR value is a
              lower bound if spin diffusion suppresses the measured torque \\
        Heavy-metal resistivity
            & $\rho_{\mathrm{HM}}=200\times10^{-8}\ \Omega\,\mathrm{m}$
            & Representative Ta/W value from the compact model;
              $\rho_{\mathrm{HM}}=200\ \mu\Omega\,\mathrm{cm}$ \\
        \hline
    \end{tabular}
    \caption{Representative properties of Pt, $\beta$-Ta, and $\beta$-W
    heavy metals~\cite{sot-stt-compact-model,liu2012ta-spin-hall,
    pai2012w-spin-hall}.}
    \label{tab:heavy-metal-properties}
\end{table}

These values are not universal elemental constants.  In particular, Pai et al.
find that W changes strongly with film phase and resistivity: a mixed
$\alpha$/$\beta$ film gives $\lvert\theta_{\mathrm{SH}}\rvert=0.18\pm0.02$,
whereas their predominantly $\alpha$-W film gives only the upper bound
$\lvert\theta_{\mathrm{SH}}\rvert<0.07$~\cite{pai2012w-spin-hall}.  The
material-specific tables below omit $\theta_{\mathrm{SH}}$; select a value from
the heavy-metal source and sign convention used by the model.

The compact-model README compares reported resistivities of
$190\ \mu\Omega\,\mathrm{cm}$ for Ta and $170$--$260\ \mu\Omega\,\mathrm{cm}$
for W, then uses the representative value
$\rho_{\mathrm{HM}}=200\ \mu\Omega\,\mathrm{cm}$~\cite{sot-stt-compact-model}.

\subsection{Ferromagnet: CoFeB}

The CoFeB free layer is represented by a one-sublattice macrospin LLGS model
with perpendicular anisotropy.  The anisotropy field and saturation magnetization
below are the defaults of the SOT/STT compact model; the remaining magnetic and
drive parameters come from the model of Zhang et
al.~\cite{sot-stt-compact-model,zhang2024revisiting}.  The compact-model
parameter named \texttt{Hk} is a flux-density quantity in tesla.  Its
corresponding $\vect{H}$ field in $\mathrm{A\,m^{-1}}$ is obtained by dividing
by $\mu_0$.

\begin{table}[htbp]
    \centering
    \small
    \renewcommand{\arraystretch}{1.15}
    \begin{tabular}{p{0.25\linewidth}p{0.28\linewidth}p{0.36\linewidth}}
        \hline
        Quantity & Value & Role or convention \\
        \hline
        Uniaxial anisotropy field
            & $B_k=1.5\ \mathrm{T}$
            & Compact-model parameter \texttt{Hk};
              $H_k\approx1.19\times10^6\ \mathrm{A\,m^{-1}}$ \\
        Saturation magnetization
            & $M_s=1000\ \mathrm{emu\,cm^{-3}}$
            & $M_s=1.00\times10^6\ \mathrm{A\,m^{-1}}$;
              $\mu_0M_s\approx1.26\ \mathrm{T}$ \\
        Gilbert damping
            & $\alpha=0.02$
            & Magnetic model parameter \\
        \hline
    \end{tabular}
    \caption{Selected CoFeB macrospin parameters from the SOT/STT compact model
    and Zhang et al.~\cite{sot-stt-compact-model,zhang2024revisiting}.}
    \label{tab:cofeb-parameters}
\end{table}

\paragraph{Parameter justification.}
Based on Nat. Mater. \textbf{9}, 721 (2010).,
$H_{\mathrm{an}}=2(K_{\mathrm{bulk}}+K_i/t_{\mathrm{FL}})/M_s$ with
$K_{\mathrm{bulk}}=2.245\times10^5\ \mathrm{J\,m^{-3}}$,
$K_i=1.286\times10^{-3}\ \mathrm{J\,m^{-2}}$,
$M_s=1.58\ \mathrm{T}=1257\ \mathrm{emu\,cm^{-3}}$, when we use
$t_{\mathrm{FL}}=1.2\ \mathrm{nm}$, it gives $H_{\mathrm{an}}=2.06\ \mathrm{T}$.
Here we approximate $M_s=1000\ \mathrm{emu\,cm^{-3}}$,
$H_k=1.5\ \mathrm{T}$~\cite{sot-stt-compact-model}.

\subsection{Ferrimagnet: GdFeCo}

Zhang et al. use a random two-dimensional atomistic
$(\mathrm{FeCo})_{1-x}\mathrm{Gd}_x$ lattice driven by damping-like STT,
whereas Zhu et al. use a one-dimensional atomistic GdFeCo model driven by
SOT~\cite{zhang2022spatially,zhu2020terahertz}.  Both are fixed-length LLG or
LLGS models.  The 2022 paper states that its parameters are otherwise the same
as those in the 2020 paper, but replaces the single exchange coefficient with
three bond-specific values.  Table~\ref{tab:gdfeco-parameters} therefore keeps
the two published exchange parameterizations separately rather than selecting
or averaging them.

\begin{table}[htbp]
    \centering
    \small
    \renewcommand{\arraystretch}{1.15}
    \begin{tabular}{p{0.19\linewidth}p{0.31\linewidth}p{0.31\linewidth}}
        \hline
        Quantity & Zhu et al. (2020) & Zhang et al. (2022) \\
        \hline
        Spatial model
            & 1D atomistic lattice with lattice constant $d=0.4\ \mathrm{nm}$;
              strip length and width are $400$ and $40\ \mathrm{nm}$
            & Random 2D square lattice; each interior site has four nearest
              neighbors; the main phase diagrams use 100 atoms \\
        Exchange
            & $A_{\mathrm{ex}}=+3\ \mathrm{meV}
              =+4.8065\times10^{-22}\ \mathrm{J}$
            & $A_{\mathrm{Gd-Gd}}=-1.26\times10^{-21}\ \mathrm{J}$;
              $A_{\mathrm{FeCo-FeCo}}=-2.83\times10^{-21}\ \mathrm{J}$;
              $A_{\mathrm{FeCo-Gd}}=+1.09\times10^{-21}\ \mathrm{J}$ \\
        Perpendicular anisotropy
            & $K_{\mathrm{TM}}=K_{\mathrm{RE}}=0.04\ \mathrm{meV}$
            & same as Zhu et al. (2020) \\
        Hard-axis anisotropy
            & $\kappa_{\mathrm{TM}}=\kappa_{\mathrm{RE}}=0.2\ \mu\mathrm{eV}$
            & Not included in the stated Hamiltonian \\
        DMI
            & $D=0.128\ \mathrm{meV}$
            & Not included in the stated Hamiltonian \\
        Gilbert damping
            & $\alpha_{\mathrm{TM}}=\alpha_{\mathrm{RE}}=0.015$
            & same as Zhu et al. (2020) \\
        Land\'e factors
            & $g_{\mathrm{TM}}=2.2$; $g_{\mathrm{RE}}=2.0$
            & same as Zhu et al. (2020) \\
        Sublattice magnetizations
            & $M_{s,\mathrm{TM}}=1149\ \mathrm{kA\,m^{-1}}$;
              $M_{s,\mathrm{RE}}=1012\ \mathrm{kA\,m^{-1}}$ at
              $T=300\ \mathrm{K}$
            & same as Zhu et al. (2020) \\
        Ferrimagnet thickness
            & $t_{\mathrm{FiM}}=0.4\ \mathrm{nm}$
            & same as Zhu et al. (2020) \\
        Drive parameter
            & Spin-Hall angle $\theta_{\mathrm{SH}}=0.2$; field-like to
              damping-like ratio $\beta$ is varied rather than fixed in the
              baseline list
            & STT efficiency $\eta$ is introduced but no numerical value is
              stated \\
        Compensation composition
            & Not stated as a composition in the parameter block
            & $x_{\mathrm{MC}}=0.23$;
              $x_{\mathrm{AMC}}=0.21$ \\
        \hline
    \end{tabular}
    \caption{Atomistic GdFeCo parameter sets retained in their native model
    conventions~\cite{zhu2020terahertz,zhang2022spatially}.}
    \label{tab:gdfeco-parameters}
\end{table}

The exchange values are deliberately not converted into a single common
coefficient.  The 2020 value belongs to a one-dimensional nearest-neighbor
Hamiltonian, while the 2022 values distinguish the chemical identity of each
bond in a random two-dimensional alloy.  Their lattice topology and bond
statistics must be retained with the numerical values when either set is used.

\subsection{Noncollinear Antiferromagnet:
\texorpdfstring{Mn$_3$Sn}{Mn3Sn}}

The three-sublattice macrospin switching model of Xu et al. and the later
spatially resolved atomistic domain-wall model use the same magnetic moment,
DMI, anisotropy, damping, and spin-Hall angle, but different exchange
coefficients~\cite{xu2024mn3sn,xu2025mn3sn-domain-wall}.  The values must
therefore remain attached to their model and interaction convention.  Both
Hamiltonians take positive $A$ to favor antiparallel alignment, and the
crystalline easy axes point toward the nearest Sn atoms.  The atomistic
domain-wall paper additionally states that its positive $D$ favors clockwise
order.

\begin{table}[htbp]
    \centering
    \small
    \renewcommand{\arraystretch}{1.15}
    \begin{tabular}{p{0.23\linewidth}p{0.29\linewidth}p{0.38\linewidth}}
        \hline
        Quantity & Value & Model scope or convention \\
        \hline
        Mn moment
            & $m_s=3\mu_B$
            & Per Mn site in both models \\
        Nearest-neighbor exchange
            & $A=17.53\ \mathrm{meV}$
            & Three-sublattice macrospin switching model \\
        Nearest-neighbor exchange
            & $A=9\ \mathrm{meV}$
            & Spatial atomistic domain-wall model \\
        DMI coefficient
            & $D=0.833\ \mathrm{meV}$
            & Same-layer interaction in both models \\
        Crystalline anisotropy
            & $K=0.196\ \mathrm{meV}$
            & Single-site Hamiltonian coefficient in both models \\
        Gilbert damping
            & $\alpha=0.003$
            & Same in both fixed-length models \\
        Unit-cell volume
            & $V_{\mathrm{cell}}=0.3778\ \mathrm{nm^3}$
            & Macrospin model; $M_s=6m_s/V_{\mathrm{cell}}
              \approx4.42\times10^5\ \mathrm{A\,m^{-1}}$ \\
        Mn$_3$Sn thickness
            & $d=30\ \mathrm{nm}$
            & Macrospin SOT model geometry \\
        Time step
            & $\Delta t=5\ \mathrm{fs}$
            & Fourth-order Runge--Kutta macrospin integration \\
        Nanoribbon length
            & $5.7\,\mu\mathrm{m}$
            & Spatial atomistic domain-wall model geometry \\
        \hline
    \end{tabular}
    \caption{Mn$_3$Sn parameter sets used in the macrospin switching and
    atomistic domain-wall models~\cite{xu2024mn3sn,xu2025mn3sn-domain-wall}.}
    \label{tab:mn3sn-parameters}
\end{table}

The Mn$_3$Sn entries are not by themselves a complete crystallographic
simulation specification: a reproducible atomistic calculation must also retain
the source model's lattice, site mapping, bond list, DMI-vector orientations,
boundary conditions, and pair-counting convention.
